\centerline{\bf \Huge Математический бой. Старшая группа}

\begin{flushright}
        \scriptsize{Только неудачники просят о смерти, проиграв бой... \\Вместо того, чтобы принять поражение, после которого ты должен был умереть!}
\end{flushright}

\problem{1}
 На летних сборах в ЭлМШ Герман решил не решать геометрию вовсе. Баир же решил решать только геометрию. Вскоре РВ посмотрел на рейтинговую таблицу и заметил, что модуль разности их баллов кратен 11. Какое количество задач они не могли решить в сумму даже в теории. (приведите все такие варианты и объясните, почему других нет.)

\problem{2}
 Действительные числа $a, b, c, d$ таковы, что $a+b = cd$ и $c+d = ab.$ Докажите неравенство $(a+1)(b+1)(c+1)(d+1) \geq 0$.

\problem{3}
Дан квадрат со стороной, большей $1000$. Докажите, что из него можно вырезать не более $4$ клеток таким образом, что остаток можно разрезать на прямоугольники $1\times4$.

\problem{4}
Дан треугольник ABC, где $\angle A = 70^{\circ}, \angle B = 30^{\circ}$ и $\angle C = 80^{\circ}$. Точка $D$ внутри треугольника такова, что $\angle DAB = 30^{\circ}$ и $\angle DBA = 10^{\circ}.$ Найдите угол $\angle DCB.$

\problem{5}
На сборах ЭлМШ $100$ учеников, каждый из которых может быть Киви, Лососем или Шиншиллой (Шиншиллы говорят правду, Лососи всегда врут, а Киви могут говорить как правду, так и ложь). $50$ участников сборов сказали: «На сборах Шиншилл больше, чем Киви», а $50$ остальных: «На сборах Лососей больше, чем Киви». Докажите, что Киви на сборах не меньше $25.$


\problem{6}
Будем называть точку плоскости узлом, если обе её координаты --- целые числа. Внутри некоторого треугольника с вершинами в узлах лежит ровно два узла (возможно, какие-то еще узлы лежат на его сторонах). Докажите, что прямая, проходящая через эти два узла, либо проходит через одну из вершин треугольника, либо параллельна одной из его сторон.

%https://problems.ru/view_problem_details_new.php?id=32889


\problem{7}
6 идущих подряд натуральных чисел каким-то образом разбивают на две тройки чисел. Затем считают произведения чисел в этих тройках и из большего произведения вычитают меньшее. Какое наименьшее значение может принимать результат?

\problem{8}
В треугольнике ABC проведена биссектриса $AL$. Точка $P$ --- основание перпендикуляра, опущенного из точки $B$ на отрезок $AL$. Из точки P на сторону $AB$ опущена высота $PQ$. На стороне AB выбрана такая точка $R$, что $RQ = AC/4.$ Докажите, что $AB+BC \geq 4PR.$